\documentclass[11pt,a4paper]{article}

% --- Pacchetti Matematici Fondamentali ---
\usepackage{amsmath}   % Strutture matematiche avanzate (es. align)
\usepackage{amssymb}   % Simboli matematici e font speciali (es. \mathbb pour i reali R)
\usepackage{amsfonts}  % Font matematici aggiuntivi
\usepackage{amsthm}    % Gestione di teoremi, definizioni e dimostrazioni

% --- Configurazione dei Teoremi ---
\newtheorem{theorem}{Teorema}
\newtheorem{lemma}[theorem]{Lemma}
\newtheorem{proposition}[theorem]{Proposizione}
\theoremstyle{definition}
\newtheorem{definition}[theorem]{Definizione}
\newtheorem{example}[theorem]{Esempio}
\theoremstyle{remark}
\newtheorem{remark}[theorem]{Nota}

\begin{document}



\begin{theorem}

	Suppose $\mathbb{E}(exp(i \xi \Gamma)) = exp(i m \xi - \frac{\sigma^2 \xi^2}{2})$. Then $m = \mathbb{E}(\Gamma)$ and $\sigma^2 = Var(\Gamma)$.

\end{theorem}




\begin{proof}

	We only examine the case for $m$, since the case for $\sigma^2$ is similar.
	Observe that $\mathbb{E}(\frac{d}{d\xi} exp(i \xi \Gamma)) = \frac{d}{d\xi} \mathbb{E} exp(i \xi \Gamma)$ (why??).
	Evaluating both sides we get:
	\[
		\mathbb{E}(i \Gamma \cdot \exp(i \xi \Gamma)) = (i m - \sigma^2 \xi) \cdot exp(i m \xi - \frac{\sigma^2 \xi^2}{2})
	\]

	Evaluating in $\xi = 0$ we obtain:

	\[
		i \cdot \mathbb{E}(\Gamma) = i \cdot m
	\]


\end{proof}


\end{document}

